\section{Conclusion}\label{sec:conclusion}

We designed and evaluated a multimodal affect recognition component for streaming video information systems over six design cycles. Four requirements guided the design: accuracy for unseen people, a compute budget, label-free handling of degradations, and accountability. On a subject-independent evaluation of all 1440 audio-visual speech clips in RAVDESS, late fusion of self-supervised speech and vision--language features reached 82.6\% accuracy over eight emotions. A cascade that runs the speech model first cut compute by 48\%, or by 21\% when both channels were monitored. The evaluation also showed a failure that clean-data benchmarks hide. Under audio noise, clean-trained models became more confident while they became less accurate, which disabled both the cascade and confidence-based drift monitoring without warning. Per-modality monitors on the inputs avoided this. Gating on their stream-level alarms protected clean-trained models. Training on degraded data protected the models more, but it turned the same gating into a liability, so gating should respond only to shifts outside the training envelope. A confirmatory cycle, with the design of Cycles~1--4 frozen, replicated these findings on 91 further speakers from CREMA-D. It also showed that monitors should describe the quality of the signal rather than its content. Quality-descriptor monitors detected every degradation in our test streams, including recorded noise and H.264 coding, and raised fewer false alarms on speaker changes. A change of corpus remained the largest threat: it cut accuracy by 31--39 points while confidence stayed high, and monitors fitted on the acted corpora, unlike confidence, noticed it. A final cycle on spontaneous speech from 1254 online videos of CMU-MOSEI confirmed the gain from fusion, the cascade's saving, the silent failure and the value of robust training at a much lower level of accuracy (55.8\% over four classes). It also qualified the monitoring results. On online video, whose clean data already vary widely, the monitors raised almost no false alarms, but the quality monitors missed recorded noise, compression and video coding; apart from compression, these did the models little harm. Models trained on acted data were confidently wrong on online video, whereas a model trained on online video stayed calibrated, and well above chance, on acted data. Our results suggest that a monitor is only as sensitive as its reference data are narrow, which makes the choice of reference data a design decision. We expect these interactions to hold for other audio-visual affect pipelines that fuse separately trained channels, but they need to be confirmed on degradations captured in the field, with the language channel, and with the operators who act on the alarms.
